%==============================================================================
%  lesson02.tex
%  Well-posedness of the linear stationary micropolar system
%  Repository: micropolar-fem
%  Author: M. Javed Bukhari
%
%  Requires macros.tex in the same folder.
%  Compile with: pdflatex lesson02.tex   (run it twice, for the contents page)
%==============================================================================

\documentclass[11pt,a4paper]{article}

% Shared macros for the theory notes.
% Include with % Shared macros for the theory notes.
% Include with % Shared macros for the theory notes.
% Include with \input{macros} after \begin{document} preamble packages.

\newcommand{\Om}{\Omega}
\newcommand{\dOm}{\partial\Omega}
\newcommand{\dx}{\,\mathrm{d}x}
\newcommand{\ds}{\,\mathrm{d}s}

\newcommand{\Ltwo}{L^2(\Om)}
\newcommand{\Ltwozero}{L^2_0(\Om)}
\newcommand{\Hone}{H^1(\Om)}
\newcommand{\Honezero}{H^1_0(\Om)}

\newcommand{\normL}[1]{\left\lVert #1 \right\rVert_{L^2(\Om)}}
\newcommand{\normH}[1]{\left\lVert #1 \right\rVert_{H^1(\Om)}}
\newcommand{\ip}[2]{\left( #1 , #2 \right)}

\newcommand{\grad}{\nabla}
\newcommand{\lap}{\Delta}
\newcommand{\dn}[1]{\frac{\partial #1}{\partial n}}

\newtheorem{theorem}{Theorem}[section]
\newtheorem{lemma}[theorem]{Lemma}
\newtheorem{proposition}[theorem]{Proposition}
\theoremstyle{definition}
\newtheorem{definition}[theorem]{Definition}
\newtheorem{remark}[theorem]{Remark}
 after \begin{document} preamble packages.

\newcommand{\Om}{\Omega}
\newcommand{\dOm}{\partial\Omega}
\newcommand{\dx}{\,\mathrm{d}x}
\newcommand{\ds}{\,\mathrm{d}s}

\newcommand{\Ltwo}{L^2(\Om)}
\newcommand{\Ltwozero}{L^2_0(\Om)}
\newcommand{\Hone}{H^1(\Om)}
\newcommand{\Honezero}{H^1_0(\Om)}

\newcommand{\normL}[1]{\left\lVert #1 \right\rVert_{L^2(\Om)}}
\newcommand{\normH}[1]{\left\lVert #1 \right\rVert_{H^1(\Om)}}
\newcommand{\ip}[2]{\left( #1 , #2 \right)}

\newcommand{\grad}{\nabla}
\newcommand{\lap}{\Delta}
\newcommand{\dn}[1]{\frac{\partial #1}{\partial n}}

\newtheorem{theorem}{Theorem}[section]
\newtheorem{lemma}[theorem]{Lemma}
\newtheorem{proposition}[theorem]{Proposition}
\theoremstyle{definition}
\newtheorem{definition}[theorem]{Definition}
\newtheorem{remark}[theorem]{Remark}
 after \begin{document} preamble packages.

\newcommand{\Om}{\Omega}
\newcommand{\dOm}{\partial\Omega}
\newcommand{\dx}{\,\mathrm{d}x}
\newcommand{\ds}{\,\mathrm{d}s}

\newcommand{\Ltwo}{L^2(\Om)}
\newcommand{\Ltwozero}{L^2_0(\Om)}
\newcommand{\Hone}{H^1(\Om)}
\newcommand{\Honezero}{H^1_0(\Om)}

\newcommand{\normL}[1]{\left\lVert #1 \right\rVert_{L^2(\Om)}}
\newcommand{\normH}[1]{\left\lVert #1 \right\rVert_{H^1(\Om)}}
\newcommand{\ip}[2]{\left( #1 , #2 \right)}

\newcommand{\grad}{\nabla}
\newcommand{\lap}{\Delta}
\newcommand{\dn}[1]{\frac{\partial #1}{\partial n}}

\newtheorem{theorem}{Theorem}[section]
\newtheorem{lemma}[theorem]{Lemma}
\newtheorem{proposition}[theorem]{Proposition}
\theoremstyle{definition}
\newtheorem{definition}[theorem]{Definition}
\newtheorem{remark}[theorem]{Remark}


\title{Well-Posedness of the Linear Stationary Micropolar System\\
on the Divergence-Free Space}
\author{M. Javed Bukhari}
\date{\today}

\begin{document}
\maketitle

\begin{abstract}
In Lesson 1 I worked through three scalar elliptic problems and showed that
boundedness and coercivity of the bilinear form are what give well-posedness.
These notes carry the same programme through for a coupled system, namely the
two dimensional stationary linear micropolar equations, restricted to
divergence free velocity fields so that the pressure disappears from the
formulation. The new difficulty is that the coupling terms have no sign, so
they can work against coercivity, and the tool for dealing with them is Young's
inequality with a free parameter. The main result is that coercivity holds with
a constant that does not involve the vortex viscosity at all, while
boundedness gives a constant that grows linearly in it. I explain at the end
what that difference means for the numerical behaviour of the system, and why
it is what one should expect physically.
\end{abstract}

\tableofcontents

%==============================================================================
\section{Results we will use}
%==============================================================================

Everything in these notes is built from the results collected here. I have
stated them all in one place because the same few are used again and again, and
it is easier to follow the argument if it is clear at every step which one is
being applied. The first three are from Lesson 1. Lemmas
\ref{lem:young}, \ref{lem:rot} and \ref{lem:curldiv} are new.

Throughout, $\Om \subset \R^{2}$ is a bounded Lipschitz domain with boundary
$\dOm$, all derivatives are weak derivatives, and I write
$\partial_{1} = \partial/\partial x$ and $\partial_{2} = \partial/\partial y$.
The $L^{2}$ inner product is
\[
  \ip{\phi}{\psi} = \int_\Om \phi \psi \dx
  \qquad \text{or} \qquad
  \ip{\boldsymbol\phi}{\boldsymbol\psi} = \int_\Om \boldsymbol\phi \cdot \boldsymbol\psi \dx
\]
depending on whether the arguments are scalars or vectors.

\begin{lemma}[Cauchy--Schwarz]\label{lem:cs}
For $\phi, \psi \in \Ltwo$,
\[
  \abs{\ip{\phi}{\psi}} \;\le\; \normL{\phi}\,\normL{\psi}.
\]
\end{lemma}

\begin{lemma}[Poincar\'e--Friedrichs]\label{lem:poincare}
There is a constant $C_\Om > 0$ depending only on $\Om$ such that
\[
  \normL{\phi} \;\le\; C_\Om \normL{\grad \phi}
  \qquad \text{for every } \phi \in \Honezero .
\]
For a vector field $\bv \in [\Honezero]^{2}$ the same bound holds with the same
constant, by applying it to each component and adding.
\end{lemma}

\begin{lemma}[Lax--Milgram]\label{lem:laxmilgram}
Let $V$ be a Hilbert space and let $a : V \times V \to \R$ be bilinear with
\[
  \abs{a(u,v)} \le M \lVert u \rVert_{V} \lVert v \rVert_{V}
  \qquad \text{and} \qquad
  a(u,u) \ge \alpha \lVert u \rVert_{V}^{2}
\]
for some constants $M, \alpha > 0$. Let $F \in V'$ be a bounded linear
functional. Then there is exactly one $u \in V$ with $a(u,v) = F(v)$ for all
$v \in V$, and it satisfies $\lVert u \rVert_{V} \le \alpha^{-1} \lVert F
\rVert_{V'}$.
\end{lemma}

\begin{lemma}[Young's inequality with a free parameter]\label{lem:young}
For $a, b \in \R$ and any $\varepsilon > 0$,
\[
  \abs{ab} \;\le\; \frac{\varepsilon}{2} a^{2} + \frac{1}{2\varepsilon} b^{2}.
\]
\end{lemma}

\begin{proof}
Since squares are non-negative,
\[
  \left( \sqrt{\varepsilon}\, a - \frac{b}{\sqrt{\varepsilon}} \right)^{2} \ge 0 .
\]
Expanding the square,
\[
  \varepsilon a^{2} - 2ab + \frac{b^{2}}{\varepsilon} \ge 0 ,
\]
so $2ab \le \varepsilon a^{2} + \varepsilon^{-1} b^{2}$, and dividing by $2$
gives the result for $ab$. Replacing $b$ by $-b$ gives the same bound for
$-ab$, so the inequality holds for $\abs{ab}$.
\end{proof}

\begin{remark}[What the free parameter is for]\label{rem:whyeps}
This is the tool I will use to handle the coupling terms, so it is worth being
clear about what it does. Suppose an estimate produces a cross term $ab$ that we
cannot control directly, but we do have squared terms $a^{2}$ and $b^{2}$
available with known positive coefficients. Young's inequality converts the
cross term into two squares, and the parameter $\varepsilon$ decides how the
cost is split between them. Taking $\varepsilon$ small puts most of the cost on
$b^{2}$ and very little on $a^{2}$. Taking $\varepsilon$ large does the
opposite. So we choose $\varepsilon$ so that each piece is small enough to be
swallowed by a term we already have. This is what is meant by absorbing a term,
and it is the single most used move in this subject.
\end{remark}

For the next two results I use the following formula, which is just one
dimensional integration by parts applied in the $x_{i}$ direction. For scalar
functions $\phi, \psi$ and either index $i \in \{1,2\}$,
\begin{equation}\label{eq:master}
  \int_\Om (\partial_{i}\phi)\psi \dx
  = \int_{\dOm} \phi \psi\, n_{i} \ds - \int_\Om \phi (\partial_{i}\psi) \dx ,
\end{equation}
where $n_{i}$ is the $i$th component of the outward unit normal. If either
$\phi$ or $\psi$ vanishes on $\dOm$ the boundary term disappears and moving a
derivative from one factor to the other costs a minus sign. Every identity
below comes from \eqref{eq:master} applied one index at a time.

In two dimensions we use the operators
\[
  \grad w = (\partial_{1}w,\ \partial_{2}w),
  \qquad
  \dive \bv = \partial_{1}v_{1} + \partial_{2}v_{2},
\]
\[
  \curl \bv = \partial_{1}v_{2} - \partial_{2}v_{1},
  \qquad
  \grad^{\perp} w = (\partial_{2}w,\ -\partial_{1}w).
\]
The last two are special to two dimensions. In three dimensions the curl of a
vector is a vector, but in two dimensions it has only one non-zero component so
we treat it as a scalar. The operator $\grad^{\perp}$ is the gradient rotated
by ninety degrees, and it turns a scalar into a vector. These two operators are
how the micropolar system converts between the vector velocity and the scalar
microrotation.

\begin{lemma}[Rotation identity]\label{lem:rot}
For $w \in \Honezero$ and $\bv \in [\Honezero]^{2}$,
\[
  \ip{\grad^{\perp} w}{\bv} = \ip{w}{\curl \bv}.
\]
\end{lemma}

\begin{proof}
Expanding the dot product into components,
\[
  \ip{\grad^{\perp} w}{\bv}
  = \int_\Om \grad^{\perp} w \cdot \bv \dx
  = \int_\Om (\partial_{2}w) v_{1} \dx - \int_\Om (\partial_{1}w) v_{2} \dx .
\]
Apply \eqref{eq:master} to each term with $\phi = w$, moving the derivative off
$w$ and onto the component of $\bv$. The boundary terms vanish because $v_{1}$
and $v_{2}$ have zero trace. This gives
\[
  \int_\Om (\partial_{2}w) v_{1} \dx = - \int_\Om w (\partial_{2}v_{1}) \dx,
  \qquad
  - \int_\Om (\partial_{1}w) v_{2} \dx = + \int_\Om w (\partial_{1}v_{2}) \dx .
\]
Adding,
\[
  \ip{\grad^{\perp} w}{\bv}
  = \int_\Om w \left( \partial_{1}v_{2} - \partial_{2}v_{1} \right) \dx
  = \int_\Om w \,\curl \bv \dx
  = \ip{w}{\curl \bv}. \qedhere
\]
\end{proof}

\begin{lemma}[Curl and divergence decomposition]\label{lem:curldiv}
For $\bv \in [\Honezero]^{2}$,
\[
  \normL{\grad \bv}^{2} = \normL{\curl \bv}^{2} + \normL{\dive \bv}^{2} ,
\]
and consequently
\[
  \normL{\curl \bv} \le \normL{\grad \bv} .
\]
\end{lemma}

\begin{proof}
Write out all three quantities in components:
\[
  \normL{\grad \bv}^{2}
  = \int_\Om \left[ (\partial_{1}v_{1})^{2} + (\partial_{2}v_{1})^{2}
  + (\partial_{1}v_{2})^{2} + (\partial_{2}v_{2})^{2} \right] \dx ,
\]
\[
  \normL{\curl \bv}^{2}
  = \int_\Om \left[ (\partial_{1}v_{2})^{2}
  - 2(\partial_{1}v_{2})(\partial_{2}v_{1}) + (\partial_{2}v_{1})^{2} \right] \dx ,
\]
\[
  \normL{\dive \bv}^{2}
  = \int_\Om \left[ (\partial_{1}v_{1})^{2}
  + 2(\partial_{1}v_{1})(\partial_{2}v_{2}) + (\partial_{2}v_{2})^{2} \right] \dx .
\]
Adding the last two, the four squared terms are exactly those appearing in
$\normL{\grad \bv}^{2}$, and what is left over is
\begin{equation}\label{eq:leftover}
  2\int_\Om (\partial_{1}v_{1})(\partial_{2}v_{2}) \dx
  - 2\int_\Om (\partial_{1}v_{2})(\partial_{2}v_{1}) \dx .
\end{equation}
So the claim will follow once we show that the two integrals in
\eqref{eq:leftover} are equal.

Start from the first one and apply \eqref{eq:master} with $\phi = v_{2}$,
$\psi = \partial_{1}v_{1}$ and $i = 2$, that is, move $\partial_{2}$ off
$v_{2}$. The boundary term vanishes because $v_{2}$ has zero trace, so
\[
  \int_\Om (\partial_{1}v_{1})(\partial_{2}v_{2}) \dx
  = - \int_\Om (\partial_{2}\partial_{1}v_{1}) v_{2} \dx .
\]
Since mixed partial derivatives commute we may write
$\partial_{2}\partial_{1}v_{1} = \partial_{1}(\partial_{2}v_{1})$, so the
integrand again has the form $(\partial_{1}\phi)\psi$ with
$\phi = \partial_{2}v_{1}$ and $\psi = v_{2}$. Applying \eqref{eq:master} a
second time, now with $i = 1$, and noting the boundary term vanishes for the
same reason,
\[
  - \int_\Om \partial_{1}(\partial_{2}v_{1}) v_{2} \dx
  = + \int_\Om (\partial_{2}v_{1})(\partial_{1}v_{2}) \dx .
\]
Chaining the two steps,
\[
  \int_\Om (\partial_{1}v_{1})(\partial_{2}v_{2}) \dx
  = \int_\Om (\partial_{2}v_{1})(\partial_{1}v_{2}) \dx ,
\]
so \eqref{eq:leftover} is zero and the identity follows. Two applications of
\eqref{eq:master} produce two minus signs, which cancel, and that is the whole
proof.

For the inequality, $\normL{\dive \bv}^{2} \ge 0$ because it is a squared norm,
so dropping it from the right hand side can only make that side smaller:
\[
  \normL{\grad \bv}^{2} \ge \normL{\curl \bv}^{2} ,
\]
and taking square roots gives the claim.
\end{proof}

\begin{remark}
The move at the end, discarding a non-negative summand to get an inequality, is
the same one used in Lesson 1 to obtain $\normL{\grad \phi} \le \normH{\phi}$
from $\normH{\phi}^{2} = \normL{\phi}^{2} + \normL{\grad \phi}^{2}$. It comes up
constantly and is worth recognising on sight.
\end{remark}

%==============================================================================
\section{The equations and what the terms mean}
%==============================================================================

\subsection{The model}

The micropolar model, due to Eringen, describes a fluid whose suspended
microstructure carries its own rotational degrees of freedom. Ordinary
Navier--Stokes treats fluid particles as points with a velocity and nothing
else. The micropolar model gives them a spin as well, and tracks it as a second
field. It is used for blood, where the red cells rotate, for ferrofluids, for
liquid crystals, and for lubricants with suspended additives.

In two dimensions, in the stationary linear case, the system is
\begin{align}
  -(\nu + \nu_{r}) \lap \bu + \grad p
    &= 2\nu_{r} \grad^{\perp} w + \mathbf{f} , \label{eq:mom} \\
  \dive \bu &= 0 , \label{eq:incomp} \\
  -(c_{a} + c_{d}) \lap w + 4\nu_{r} w
    &= 2\nu_{r} \curl \bu + g , \label{eq:micro}
\end{align}
with $\bu = \mathbf{0}$ and $w = 0$ on $\dOm$. The unknowns are the velocity
$\bu$, the pressure $p$, and the microrotation $w$, which is the angular
velocity of the microstructure.

\subsection{What each term is doing}

It is worth going through the terms one at a time, because the whole point of
the analysis is to find out how they interact, and that is much easier to follow
if their meanings are already clear.

\begin{itemize}

\item $-(\nu + \nu_{r}) \lap \bu$ is viscous diffusion of momentum. The
coefficient is the sum of two viscosities. The first, $\nu$, is the ordinary
dynamic viscosity, the same one as in Navier--Stokes. The second, $\nu_{r}$, is
called the vortex viscosity, and it is the extra resistance the fluid feels
because the microstructure is there at all. Note that this term is where
$\nu_{r}$ first appears, and it appears with a plus sign, so at this stage it
looks like it can only help.

\item $\grad p$ is the pressure gradient. Pressure enters the momentum equation
only through its gradient, which will matter shortly.

\item $2\nu_{r}\grad^{\perp} w$ is the first coupling term. It says that where
the microstructure is spinning, that spin drags the bulk fluid along. This is a
transfer of angular momentum from the microrotation field to the velocity field.

\item $\dive \bu = 0$ is incompressibility.

\item $-(c_{a} + c_{d}) \lap w$ is angular diffusion. Spin spreads out through
the microstructure in the same way that momentum spreads out through the fluid.
The constants $c_{a}$ and $c_{d}$ are angular viscosities, and they only ever
appear in the combination $c_{a} + c_{d}$ in the two dimensional case, so I will
treat that sum as a single positive constant.

\item $4\nu_{r} w$ is a damping or reaction term. It says that spin decays. If
the microstructure spins faster than the surrounding fluid rotates, friction
slows it down. It is the only zeroth order term in the whole system, and it will
turn out to be very important.

\item $2\nu_{r} \curl \bu$ is the second coupling term, the mirror image of the
first. Where the bulk fluid is rotating, it spins up the microstructure. This is
the transfer in the opposite direction.

\end{itemize}

The key observation, which is the theme of these notes, is that $\nu_{r}$
appears in four places: in the velocity diffusion, in the damping term, and in
both couplings. It is not four independent parameters. It is one physical
quantity, the vortex viscosity, appearing four times, and the coefficients
$2\nu_{r}$, $2\nu_{r}$ and $4\nu_{r}$ are related to each other, not free. That
relationship is what makes the analysis below work out as cleanly as it does.

\subsection{Why we work on the divergence-free space}

To keep the analysis to one difficulty at a time, I work here on the space of
divergence free velocity fields,
\[
  \mathbf{V}_{\mathrm{div}}
  = \left\{ \bv \in [\Honezero]^{2} : \dive \bv = 0 \right\} ,
  \qquad
  W = \Honezero .
\]
The reason is that when we derive the weak form, the pressure appears only in
the term
\[
  - \int_\Om p \,\dive \bv \dx ,
\]
and if $\bv$ is divergence free this vanishes identically. So the pressure drops
out of the formulation completely, and we are left with a problem in two
unknowns rather than three, on a single product space, which is exactly the
setting Lax--Milgram is stated for.

This is legitimate mathematics and it gives us a genuine well-posedness result.
It is not, however, how one discretises the problem, because building a finite
element basis for $\mathbf{V}_{\mathrm{div}}$ is hard. That difficulty is what
forces the mixed formulation and the inf-sup condition, and it is the subject of
the next set of notes. Doing it this way round means we get the coercivity
result cleanly here, and meet the pressure as a separate problem rather than
fighting both at once.

\subsection{The product space}

On $\mathbf{V}_{\mathrm{div}} \times W$ we use the norm
\begin{equation}\label{eq:prodnorm}
  \lVert (\bv, z) \rVert^{2}
  = \normH{\bv}^{2} + \normH{z}^{2}
  = \normL{\bv}^{2} + \normL{\grad \bv}^{2}
  + \normL{z}^{2} + \normL{\grad z}^{2} .
\end{equation}
The idea is simply to treat the pair $(\bu, w)$ as one unknown living in one
space, so that everything from Lesson 1 applies with no structural change. This
is the standard way to handle coupled systems and it is worth getting used to.

\subsection{From three equations to one}

The system \eqref{eq:mom} to \eqref{eq:micro} has three equations and three
unknowns, but the analysis below is carried out on a single equation with a
single bilinear form. It is worth doing that reduction slowly, because three
separate things happen and it is easy to lose track of which is which. The
incompressibility equation is absorbed into the choice of space, the pressure
term then disappears as a consequence, and the two remaining equations are added
together into one.

\paragraph{Step 1: the incompressibility equation becomes the space.}
The equation
\[
  \dive \bu = 0
\]
is a constraint on $\bu$ rather than a differential equation we need to test
against. So instead of multiplying it by a test function, we build it into the
space in which we look for $\bu$:
\[
  \mathbf{V}_{\mathrm{div}}
  = \left\{ \bv \in [\Honezero]^{2} : \dive \bv = 0 \right\} .
\]
Searching for $\bu$ in this space imposes $\dive \bu = 0$ exactly. That disposes
of one of the three equations, and two remain.

\paragraph{Step 2: the momentum equation.}
Multiply \eqref{eq:mom} by a test function $\bv \in \mathbf{V}_{\mathrm{div}}$
and integrate over $\Om$:
\[
  -(\nu+\nu_{r}) \int_\Om (\lap \bu)\cdot \bv \dx
  + \int_\Om \grad p \cdot \bv \dx
  = 2\nu_{r} \int_\Om \grad^{\perp} w \cdot \bv \dx
  + \int_\Om \mathbf{f}\cdot \bv \dx .
\]

For the Laplacian term, apply Green's identity componentwise exactly as in
Lesson 1. The boundary terms vanish because $\bv$ has zero trace, giving
\[
  -(\nu+\nu_{r}) \int_\Om (\lap \bu)\cdot \bv \dx
  = (\nu+\nu_{r}) \ip{\grad \bu}{\grad \bv} .
\]

For the pressure term, apply the master formula \eqref{eq:master} with
$\phi = p$ and $\psi = v_{i}$, once for each component. The boundary terms again
vanish because $\bv$ has zero trace, and each move costs a minus sign:
\[
  \int_\Om \grad p \cdot \bv \dx
  = \sum_{i=1}^{2} \int_\Om (\partial_{i} p) v_{i} \dx
  = - \sum_{i=1}^{2} \int_\Om p\, (\partial_{i} v_{i}) \dx
  = - \int_\Om p \,\dive \bv \dx .
\]
But $\bv \in \mathbf{V}_{\mathrm{div}}$, so $\dive \bv = 0$ and the entire
pressure term is zero.

This is the payoff from Step 1, and it is the whole reason for restricting the
space. It is worth noticing that the pressure vanished because of the divergence
of the \emph{test} function, not of the solution. The weak form of the momentum
equation is therefore
\begin{equation}\label{eq:weakmom}
  (\nu+\nu_{r}) \ip{\grad \bu}{\grad \bv}
  - 2\nu_{r} \ip{\grad^{\perp} w}{\bv}
  = \ip{\mathbf{f}}{\bv}
  \qquad \text{for all } \bv \in \mathbf{V}_{\mathrm{div}} .
\end{equation}

\paragraph{Step 3: the microrotation equation.}
Multiply \eqref{eq:micro} by a test function $z \in W = \Honezero$ and integrate
over $\Om$. Green's identity applies to the Laplacian term, with the boundary
term vanishing because $z$ has zero trace, and the zeroth order term $4\nu_{r}w$
contains no derivative so nothing needs to be done to it. This gives
\begin{equation}\label{eq:weakmicro}
  (c_{a}+c_{d}) \ip{\grad w}{\grad z}
  + 4\nu_{r} \ip{w}{z}
  - 2\nu_{r} \ip{\curl \bu}{z}
  = \ip{g}{z}
  \qquad \text{for all } z \in W .
\end{equation}

\paragraph{Step 4: adding the two equations.}
We now have two weak equations, \eqref{eq:weakmom} tested against every $\bv$
and \eqref{eq:weakmicro} tested against every $z$. I claim this is equivalent to
testing their sum against every pair $(\bv, z)$.

One direction is immediate: if \eqref{eq:weakmom} and \eqref{eq:weakmicro} both
hold then so does their sum, for every pair. The other direction is the one that
needs an argument. Suppose the summed equation holds for all
$(\bv, z) \in \mathbf{V}_{\mathrm{div}} \times W$. Choosing $z = 0$ makes every
term coming from \eqref{eq:weakmicro} vanish, and what is left is exactly
\eqref{eq:weakmom} for arbitrary $\bv$. Choosing $\bv = \mathbf{0}$ recovers
\eqref{eq:weakmicro} in the same way. So both original equations follow from the
sum, and the two formulations carry exactly the same information. Nothing is
lost by adding them, and it is the freedom to set one component of the test pair
to zero that makes the step reversible.

Adding \eqref{eq:weakmom} and \eqref{eq:weakmicro} therefore gives the weak
problem in its final form: find
$(\bu, w) \in \mathbf{V}_{\mathrm{div}} \times W$ such that
\begin{equation}\label{eq:weak}
  a\big( (\bu, w), (\bv, z) \big) = F(\bv, z)
  \qquad \text{for all } (\bv, z) \in \mathbf{V}_{\mathrm{div}} \times W ,
\end{equation}
where
\begin{multline}\label{eq:bilinear}
  a\big( (\bu, w), (\bv, z) \big)
  = (\nu + \nu_{r}) \ip{\grad \bu}{\grad \bv}
  + (c_{a}+c_{d}) \ip{\grad w}{\grad z}
  + 4\nu_{r} \ip{w}{z} \\
  {} - 2\nu_{r} \ip{\grad^{\perp} w}{\bv}
  - 2\nu_{r} \ip{\curl \bu}{z}
\end{multline}
and $F(\bv, z) = \ip{\mathbf{f}}{\bv} + \ip{g}{z}$.

The first three terms are things we have already met. The first is the Stokes
term. The second and third together are exactly the reaction diffusion problem
from Lesson 1, Problem C, with $c = c_{a} + c_{d}$ and $\kappa = 4\nu_{r}$. The
last two are the couplings, and they are the new difficulty, because they carry
a minus sign and there is nothing forcing them to be small. They can be
negative, and if they are large and negative they could in principle destroy
coercivity.

\begin{remark}[Why we bother adding the equations]
It would be perfectly possible to keep \eqref{eq:weakmom} and
\eqref{eq:weakmicro} as two separate equations, and nothing would be wrong with
that. The reason for combining them is that Lax--Milgram, Lemma
\ref{lem:laxmilgram}, is a statement about one bilinear form on one Hilbert
space. It says nothing about a pair of coupled equations posed on two different
spaces. By treating the pair $(\bu, w)$ as a single unknown living in the single
product space $\mathbf{V}_{\mathrm{div}} \times W$, every result from Lesson 1
becomes applicable without any structural change, and indeed the coercivity
proof below turns out to be the Lesson 1 argument with one extra absorption
step.

This packaging is the standard way of handling coupled systems and it is worth
having as a reflex. When faced with a system of several equations, look for the
product space that turns it into one.
\end{remark}

%==============================================================================
\section{Coercivity}
%==============================================================================

\subsection{Plan of the argument, and why $\varepsilon$ comes first}

Before doing the calculation I want to say what the shape of it is, because the
order of operations is not obvious the first time.

We set $\bv = \bu$ and $z = w$, which turns the first three terms of
\eqref{eq:bilinear} into norms immediately and leaves one term of unknown sign.
That bad term is a cross term: it is an inner product of two different
quantities, $\curl \bu$ and $w$, and there is no reason for it to be small. So
we bound it, and the only way to bound a cross term in terms of squares is
Cauchy--Schwarz followed by Young.

The point of Remark \ref{rem:whyeps} is that Young leaves us a free parameter,
and until we fix it we do not know how much damage has been done to each of the
good terms. So the parameter has to be chosen, and it has to be chosen at this
stage rather than later, because the whole question of whether coercivity holds
is the question of whether there exists a choice of $\varepsilon$ making all the
remaining coefficients non-negative. If no such $\varepsilon$ exists, the
argument fails and we would have to conclude something different about the
system. So finding $\varepsilon$ is not a technical detail postponed to the end.
It is the crux.

Once $\varepsilon$ is fixed we substitute it in, the constants become explicit
numbers, and after that the parameter never appears again. This is why
$\varepsilon$ is used in the coercivity proof and does not appear anywhere in
the boundedness proof. In boundedness there is nothing to absorb, because the
two arguments $(\bu,w)$ and $(\bv,z)$ are independent and there is no
possibility of one term cancelling another. Every term simply gets bounded on
its own and the constants are added up.

\subsection{Step 1: setting the arguments equal}

Put $\bv = \bu$ and $z = w$ in \eqref{eq:bilinear}. The first three terms become
norms straight away, since $\ip{\phi}{\phi} = \normL{\phi}^{2}$ directly from
the definition of the norm. There is no estimate at this step:
\[
  (\nu+\nu_{r})\ip{\grad \bu}{\grad \bu} = (\nu+\nu_{r})\normL{\grad \bu}^{2} ,
\]
\[
  (c_{a}+c_{d})\ip{\grad w}{\grad w} = (c_{a}+c_{d})\normL{\grad w}^{2} ,
\]
\[
  4\nu_{r}\ip{w}{w} = 4\nu_{r}\normL{w}^{2} .
\]

For the fourth term we use Lemma \ref{lem:rot} with $\bv = \bu$, together with
the symmetry of the $L^{2}$ inner product:
\[
  \ip{\grad^{\perp} w}{\bu} = \ip{w}{\curl \bu} = \ip{\curl \bu}{w} .
\]
So the fourth term becomes $-2\nu_{r}\ip{\curl \bu}{w}$, which is identical to
the fifth term. The two couplings therefore combine into a single term:
\[
  - 2\nu_{r}\ip{\curl \bu}{w} - 2\nu_{r}\ip{\curl \bu}{w}
  = - 4\nu_{r}\ip{\curl \bu}{w} .
\]

Collecting everything,
\begin{equation}\label{eq:identity}
  a\big( (\bu,w),(\bu,w) \big)
  = (\nu+\nu_{r})\normL{\grad \bu}^{2}
  + (c_{a}+c_{d})\normL{\grad w}^{2}
  + 4\nu_{r}\normL{w}^{2}
  - 4\nu_{r}\ip{\curl \bu}{w} .
\end{equation}
This is an identity, not an estimate. Three positive terms and one whose sign we
do not know.

\subsection{Step 2: bounding the coupling term}

Apply Cauchy--Schwarz (Lemma \ref{lem:cs}) to the last term of
\eqref{eq:identity}:
\begin{equation}\label{eq:cs-applied}
  \abs{4\nu_{r}\ip{\curl \bu}{w}}
  \le 4\nu_{r} \normL{\curl \bu} \normL{w} .
\end{equation}
We now have a product of two different norms, and nothing in
\eqref{eq:identity} to compare it against directly, since everything there is
squared. So apply Young (Lemma \ref{lem:young}) with
\[
  a = \normL{\curl \bu}, \qquad b = \normL{w},
\]
and a parameter $\varepsilon > 0$ still to be chosen:
\begin{equation}\label{eq:young-applied}
  \normL{\curl \bu}\normL{w}
  \le \frac{\varepsilon}{2}\normL{\curl \bu}^{2}
  + \frac{1}{2\varepsilon}\normL{w}^{2} .
\end{equation}

The first term on the right involves $\normL{\curl \bu}^{2}$, which does not
appear anywhere in \eqref{eq:identity}. What does appear is
$\normL{\grad \bu}^{2}$. Lemma \ref{lem:curldiv} is the bridge between them:
\begin{equation}\label{eq:curl-bound}
  \normL{\curl \bu}^{2} \le \normL{\grad \bu}^{2} .
\end{equation}
Note that on $\mathbf{V}_{\mathrm{div}}$ we actually have equality here, because
$\dive \bu = 0$, so nothing is being wasted. I write it as an inequality because
that is the form that survives if the divergence free constraint is later
relaxed.

Combining \eqref{eq:cs-applied}, \eqref{eq:young-applied} and
\eqref{eq:curl-bound},
\begin{equation}\label{eq:coupling-bound}
  \abs{4\nu_{r}\ip{\curl \bu}{w}}
  \le 4\nu_{r}\left[
    \frac{\varepsilon}{2}\normL{\grad \bu}^{2}
    + \frac{1}{2\varepsilon}\normL{w}^{2}
  \right]
  = 2\nu_{r}\varepsilon \normL{\grad \bu}^{2}
  + \frac{2\nu_{r}}{\varepsilon}\normL{w}^{2} .
\end{equation}

\subsection{Step 3: substituting back}

Since $-X \ge -\abs{X}$ for any real $X$, we may replace the coupling term in
\eqref{eq:identity} by minus the bound \eqref{eq:coupling-bound}:
\begin{align}
  a\big( (\bu,w),(\bu,w) \big)
  &\ge (\nu+\nu_{r})\normL{\grad \bu}^{2}
  + (c_{a}+c_{d})\normL{\grad w}^{2}
  + 4\nu_{r}\normL{w}^{2} \notag \\
  &\qquad {} - 2\nu_{r}\varepsilon\normL{\grad \bu}^{2}
  - \frac{2\nu_{r}}{\varepsilon}\normL{w}^{2} \notag \\[2mm]
  &= \Big( \nu + \nu_{r} - 2\nu_{r}\varepsilon \Big) \normL{\grad \bu}^{2}
  + (c_{a}+c_{d})\normL{\grad w}^{2}
  + \Big( 4\nu_{r} - \frac{2\nu_{r}}{\varepsilon} \Big) \normL{w}^{2} .
  \label{eq:with-eps}
\end{align}

\subsection{Step 4: choosing $\varepsilon$}

Now we choose the parameter. There are two things we want.

First, we want the coefficient of $\normL{\grad \bu}^{2}$ to be free of
$\nu_{r}$, that is, we want the $\nu_{r}$ contributions to cancel:
\begin{equation}\label{eq:cond1}
  \nu_{r} - 2\nu_{r}\varepsilon = 0 .
\end{equation}
Second, we need the coefficient of $\normL{w}^{2}$ to be non-negative, since a
negative coefficient there would leave a term working against us:
\begin{equation}\label{eq:cond2}
  4\nu_{r} - \frac{2\nu_{r}}{\varepsilon} \ge 0 .
\end{equation}

That is two conditions and one unknown, so it is not obvious in advance that
both can be met. Solve \eqref{eq:cond1}. Dividing by $\nu_{r} > 0$,
\[
  1 - 2\varepsilon = 0 \qquad \Longrightarrow \qquad \varepsilon = \tfrac{1}{2} .
\]
Now check \eqref{eq:cond2} at this value. With $\varepsilon = 1/2$ we have
$2\nu_{r}/\varepsilon = 4\nu_{r}$, so
\[
  4\nu_{r} - \frac{2\nu_{r}}{\varepsilon} = 4\nu_{r} - 4\nu_{r} = 0 \ \ge 0 ,
\]
which is satisfied, though only just: the coefficient is exactly zero, not
positive.

\begin{remark}[This is not a coincidence]\label{rem:notcoincidence}
The one value of $\varepsilon$ that removes $\nu_{r}$ from the velocity
coefficient is the same value that exactly exhausts the damping term
$4\nu_{r}\normL{w}^{2}$. Both constraints bind at once, with nothing to spare.

This happens because the coefficients in the micropolar model are not
independent. The two couplings carry $2\nu_{r}$ each and the damping term
carries $4\nu_{r}$, and these are tied together by the way the model is derived
from conservation of angular momentum. The algebra above is the mathematical
shadow of that physical relationship. I say more about this in Section
\ref{sec:conclusions}.
\end{remark}

Substituting $\varepsilon = 1/2$ into \eqref{eq:with-eps}, the velocity
coefficient becomes $\nu + \nu_{r} - \nu_{r} = \nu$ and the $\normL{w}^{2}$ term
disappears entirely:
\begin{equation}\label{eq:coerc-seminorm}
  \boxed{\;
  a\big( (\bu,w),(\bu,w) \big)
  \ \ge\ \nu \normL{\grad \bu}^{2} + (c_{a}+c_{d})\normL{\grad w}^{2} .
  \;}
\end{equation}
Two things are worth noticing. The coefficient of the velocity term is $\nu$,
not $\nu + \nu_{r}$: the vortex viscosity contribution was entirely consumed in
absorbing the coupling. And $\nu_{r}$ does not appear anywhere on the right hand
side.

\subsection{Step 5: converting to the full product norm}

Estimate \eqref{eq:coerc-seminorm} controls only the gradient parts. The norm
\eqref{eq:prodnorm} also contains $\normL{\bu}^{2}$ and $\normL{w}^{2}$, and we
have neither, since the $\normL{w}^{2}$ term was used up. Both fields have zero
trace, so Poincar\'e supplies what is missing. This is the Lesson 1 argument run
twice.

By Lemma \ref{lem:poincare},
\[
  \normL{\bu}^{2} \le C_\Om^{2} \normL{\grad \bu}^{2},
  \qquad
  \normL{w}^{2} \le C_\Om^{2} \normL{\grad w}^{2},
\]
so
\[
  \normH{\bu}^{2} = \normL{\bu}^{2} + \normL{\grad \bu}^{2}
  \le \left( C_\Om^{2} + 1 \right) \normL{\grad \bu}^{2} ,
\]
and identically for $w$. Rearranging both,
\begin{equation}\label{eq:poincare-inverted}
  \normL{\grad \bu}^{2} \ge \frac{\normH{\bu}^{2}}{C_\Om^{2}+1},
  \qquad
  \normL{\grad w}^{2} \ge \frac{\normH{w}^{2}}{C_\Om^{2}+1} .
\end{equation}
Substituting \eqref{eq:poincare-inverted} into \eqref{eq:coerc-seminorm},
\[
  a\big( (\bu,w),(\bu,w) \big)
  \ge \frac{\nu}{C_\Om^{2}+1}\normH{\bu}^{2}
  + \frac{c_{a}+c_{d}}{C_\Om^{2}+1}\normH{w}^{2} .
\]
Finally, replace both coefficients by the smaller of the two, exactly as in
Lesson 1 Problem C:
\[
  a\big( (\bu,w),(\bu,w) \big)
  \ge \frac{\min(\nu,\ c_{a}+c_{d})}{C_\Om^{2}+1}
  \left( \normH{\bu}^{2} + \normH{w}^{2} \right) .
\]

We have proved the following.

\begin{theorem}[Coercivity]\label{thm:coercivity}
For all $(\bu,w) \in \mathbf{V}_{\mathrm{div}} \times W$,
\[
  a\big( (\bu,w),(\bu,w) \big)
  \ \ge\ \alpha \, \lVert (\bu,w) \rVert^{2} ,
  \qquad
  \alpha = \frac{\min(\nu,\ c_{a}+c_{d})}{C_\Om^{2}+1} .
\]
In particular $\alpha$ does not depend on the vortex viscosity $\nu_{r}$.
\end{theorem}

%==============================================================================
\section{Boundedness}
%==============================================================================

\subsection{Why there is no $\varepsilon$ here}

In the coercivity proof the two arguments were the same, which is what allowed
one term to be absorbed into another. In boundedness the two arguments
$(\bu,w)$ and $(\bv,z)$ are completely independent, so no cancellation between
terms is possible and there is nothing to absorb. Each term is bounded
separately by Cauchy--Schwarz and the constants are added. Young's inequality
plays no role.

Throughout we use, repeatedly, the fact that each individual piece is bounded by
the full product norm. From \eqref{eq:prodnorm}, all four summands are
non-negative, so discarding any of them gives, for instance,
\begin{equation}\label{eq:pieces}
  \normL{\grad \bu} \le \normH{\bu} \le \lVert (\bu,w) \rVert,
  \qquad
  \normL{w} \le \normH{w} \le \lVert (\bu,w) \rVert ,
\end{equation}
and similarly for every other piece. This is the same discard-a-non-negative-term
move as before.

\subsection{The five terms}

\paragraph{Term 1.} By Cauchy--Schwarz and then \eqref{eq:pieces},
\[
  (\nu+\nu_{r})\abs{\ip{\grad \bu}{\grad \bv}}
  \le (\nu+\nu_{r}) \normL{\grad \bu}\normL{\grad \bv}
  \le (\nu+\nu_{r}) \lVert (\bu,w) \rVert \lVert (\bv,z) \rVert .
\]

\paragraph{Term 2.} Identically,
\[
  (c_{a}+c_{d})\abs{\ip{\grad w}{\grad z}}
  \le (c_{a}+c_{d}) \normL{\grad w}\normL{\grad z}
  \le (c_{a}+c_{d}) \lVert (\bu,w) \rVert \lVert (\bv,z) \rVert .
\]

\paragraph{Term 3.} Again,
\[
  4\nu_{r}\abs{\ip{w}{z}}
  \le 4\nu_{r} \normL{w}\normL{z}
  \le 4\nu_{r} \lVert (\bu,w) \rVert \lVert (\bv,z) \rVert .
\]

\paragraph{Term 4.} Here we first use Lemma \ref{lem:rot} to move the derivative
onto $\bv$, then Cauchy--Schwarz, then Lemma \ref{lem:curldiv} to replace the
curl by the full gradient:
\[
  2\nu_{r}\abs{\ip{\grad^{\perp} w}{\bv}}
  = 2\nu_{r}\abs{\ip{w}{\curl \bv}}
  \le 2\nu_{r} \normL{w}\normL{\curl \bv}
  \le 2\nu_{r} \normL{w}\normL{\grad \bv} ,
\]
and then by \eqref{eq:pieces} this is at most
$2\nu_{r} \lVert (\bu,w) \rVert \lVert (\bv,z) \rVert$.

\paragraph{Term 5.} The same, with the roles of the two arguments exchanged:
\[
  2\nu_{r}\abs{\ip{\curl \bu}{z}}
  \le 2\nu_{r} \normL{\curl \bu}\normL{z}
  \le 2\nu_{r} \normL{\grad \bu}\normL{z}
  \le 2\nu_{r} \lVert (\bu,w) \rVert \lVert (\bv,z) \rVert .
\]

\subsection{Adding up}

By the triangle inequality applied to \eqref{eq:bilinear}, the modulus of the
whole form is at most the sum of the five bounds. Collecting the constants,
\[
  M = (\nu + \nu_{r}) + (c_{a}+c_{d}) + 4\nu_{r} + 2\nu_{r} + 2\nu_{r}
  = \nu + (c_{a}+c_{d}) + 9\nu_{r} .
\]

\begin{theorem}[Boundedness]\label{thm:boundedness}
For all $(\bu,w), (\bv,z) \in \mathbf{V}_{\mathrm{div}} \times W$,
\[
  \abs{a\big( (\bu,w),(\bv,z) \big)}
  \ \le\ M \lVert (\bu,w) \rVert \lVert (\bv,z) \rVert ,
  \qquad
  M = \nu + (c_{a}+c_{d}) + 9\nu_{r} .
\]
In particular $M$ grows linearly in $\nu_{r}$.
\end{theorem}

%==============================================================================
\section{Well-posedness}
%==============================================================================

\begin{theorem}[Well-posedness]\label{thm:wellposed}
Let $\nu > 0$, $c_{a}+c_{d} > 0$, $\nu_{r} \ge 0$, and let
$\mathbf{f} \in [\Ltwo]^{2}$ and $g \in \Ltwo$. Then the weak problem
\eqref{eq:weak} has exactly one solution
$(\bu, w) \in \mathbf{V}_{\mathrm{div}} \times W$, and it satisfies
\[
  \lVert (\bu,w) \rVert
  \le \frac{1}{\alpha} \lVert F \rVert ,
  \qquad
  \alpha = \frac{\min(\nu,\ c_{a}+c_{d})}{C_\Om^{2}+1} .
\]
\end{theorem}

\begin{proof}
The product space $\mathbf{V}_{\mathrm{div}} \times W$ is a Hilbert space with
the norm \eqref{eq:prodnorm}, being a closed subspace of a product of Hilbert
spaces. The form $a$ is bilinear by inspection, bounded by Theorem
\ref{thm:boundedness}, and coercive by Theorem \ref{thm:coercivity}. The
functional $F$ is bounded because, by Cauchy--Schwarz,
\[
  \abs{F(\bv,z)} \le \normL{\mathbf{f}}\normL{\bv} + \normL{g}\normL{z}
  \le \left( \normL{\mathbf{f}} + \normL{g} \right) \lVert (\bv,z) \rVert .
\]
Lemma \ref{lem:laxmilgram} then gives existence, uniqueness, and the stated
bound.
\end{proof}

%==============================================================================
\section{Conclusions}
\label{sec:conclusions}
%==============================================================================

\subsection{Summary of the constants}

\begin{center}
\begin{tabular}{@{}lll@{}}
\toprule
Quantity & Value & Depends on $\nu_{r}$? \\
\midrule
Coercivity constant $\alpha$
  & $\dfrac{\min(\nu,\ c_{a}+c_{d})}{C_\Om^{2}+1}$ & no \\[3mm]
Boundedness constant $M$
  & $\nu + (c_{a}+c_{d}) + 9\nu_{r}$ & yes, linearly \\[2mm]
Ratio $M/\alpha$
  & $\dfrac{\big(\nu + (c_{a}+c_{d}) + 9\nu_{r}\big)(C_\Om^{2}+1)}
           {\min(\nu,\ c_{a}+c_{d})}$ & yes, linearly \\
\bottomrule
\end{tabular}
\end{center}

\subsection{The limit of small vortex viscosity}

As $\nu_{r} \to 0$, both couplings and the damping term all vanish, since each
carries a factor $\nu_{r}$. The bilinear form reduces to
\[
  a\big( (\bu,w),(\bv,z) \big) \longrightarrow
  \nu \ip{\grad \bu}{\grad \bv} + (c_{a}+c_{d}) \ip{\grad w}{\grad z} ,
\]
which is completely decoupled. The velocity solves a divergence free Stokes
problem and the microrotation solves a Poisson problem, and neither knows about
the other. This makes sense physically: $\nu_{r}$ is the coupling. With no
vortex viscosity the microstructure spins freely and exchanges nothing with the
bulk flow.

The coercivity constant survives this limit unchanged, and each of the two
constants inside the minimum is precisely the Lesson 1 constant for the
corresponding decoupled problem. So the limit is continuous, and the analysis
degenerates to the right thing, which is a good check that no error has crept
in.

\subsection{The limit of large vortex viscosity, and what it predicts numerically}

As $\nu_{r} \to \infty$ the situation is different. Coercivity never fails:
$\alpha$ is fixed and strictly positive for every $\nu_{r}$, so by Theorem
\ref{thm:wellposed} the problem remains uniquely solvable no matter how large
the vortex viscosity is. But $M$ grows without bound, so the ratio
\[
  \frac{M}{\alpha} \sim \frac{9\nu_{r}\left( C_\Om^{2}+1 \right)}
  {\min(\nu,\ c_{a}+c_{d})}
  \qquad \text{as } \nu_{r} \to \infty
\]
grows linearly. This ratio is the condition number of the problem, and it is
exactly the constant that appears in C\'ea's lemma:
\[
  \lVert (\bu,w) - (\bu_{h},w_{h}) \rVert
  \le \frac{M}{\alpha}
  \inf_{(\bv_{h},z_{h})} \lVert (\bu,w) - (\bv_{h},z_{h}) \rVert .
\]

So the analysis makes a concrete numerical prediction. The convergence
\emph{rates} of a finite element discretisation should be unaffected by
$\nu_{r}$, because they come from the interpolation estimate on the right hand
side. But the error \emph{constants} should grow linearly in $\nu_{r}$, meaning
that for a fixed mesh the computed solution gets worse as the vortex viscosity
increases, and a finer mesh is needed to recover the same accuracy. The discrete
linear system should also become progressively harder for iterative solvers, so
Krylov iteration counts should climb and preconditioning should stop being
optional.

This is a prediction that can be tested directly, and testing it is what the
parameter study in the project plan is for. Having derived the prediction from
the analysis first, rather than simply running the code and describing what
came out, is the difference between numerical analysis and computation.

\subsection{Why $\nu_{r}$ disappears from the coercivity constant}

The vortex viscosity appears four times in the equations, and yet it leaves no
trace at all in $\alpha$. I think this is the most interesting thing in these
notes, and the reason for it is physical rather than algebraic.

Coercivity is fundamentally an energy statement. It says that the bilinear form
controls the total energy of the pair $(\bu, w)$ from below, which is the same
as saying that the system dissipates energy at a rate proportional to how much
energy it has. Now consider what $\nu_{r}$ actually does. The two coupling terms
transfer angular momentum between the bulk flow and the microstructure, one in
each direction. Transfer is not dissipation: energy moved from the velocity
field to the microrotation field is still in the system. A mechanism that only
moves energy around internally cannot contribute anything to the total, so it
must cancel out of any statement about the total, and that is exactly what the
choice $\varepsilon = 1/2$ achieves.

The energy that genuinely leaves the system is dissipated by the two
diffusivities: $\nu$, through viscous shear in the bulk fluid, and
$c_{a} + c_{d}$, through angular diffusion in the microstructure. These are the
only two mechanisms that actually destroy energy, and so these are the only two
that appear in $\alpha$. The minimum appears because the system as a whole can
only dissipate as fast as its slowest channel.

This also explains Remark \ref{rem:notcoincidence}. The two conditions
\eqref{eq:cond1} and \eqref{eq:cond2} bind simultaneously because the
coefficients $2\nu_{r}$, $2\nu_{r}$ and $4\nu_{r}$ in the model are not
independent parameters. They come from a single conservation law, and the exact
cancellation in the algebra is the trace that conservation law leaves behind. If
the model had been written with arbitrary unrelated coefficients, the two
conditions would generally have been incompatible and coercivity would have
required a smallness assumption on the coupling. The fact that it does not is a
statement about the model, not about the proof.

\subsection{What comes next}

The result proved here is genuine but incomplete, because the space
$\mathbf{V}_{\mathrm{div}}$ cannot be discretised conveniently. Constructing a
finite element basis of divergence free functions is possible but awkward, and
essentially nobody does it. The practical route is to keep the pressure as an
unknown, work on $[\Honezero]^{2} \times \Ltwozero \times \Honezero$, and impose
$\dive \bu = 0$ weakly through the pressure.

The price is that coercivity is no longer enough. The pressure enters through a
constraint rather than through a coercive term, so Lax--Milgram does not apply
and must be replaced by the Babu\v{s}ka--Brezzi theory, whose extra hypothesis
is the inf-sup condition. That is the subject of Lesson 3, and it is the real
reason mixed finite element methods exist as a separate subject.

Nothing proved here is wasted. The coercivity estimate \eqref{eq:coerc-seminorm}
is exactly the first of the two Brezzi conditions, and it carries over unchanged.
What has to be added is the second.

\end{document}
